% f02_canonical_tree variant c -- FBT(8) built in four steps: the leaves are
% placed, then each level is folded upward, and at every step the set of legal
% moves has exactly one member per surviving pair.
% Data: Definitions 1-3 of main.tex (sec:abi-tree); the heap arithmetic and the
% leaf placement are the released implementation
% (experiments/w_final/lib/pinned_fbt.py: leaf_index, children) and CBTree in
% src/sprs_core.py.  The N=8 root word is the N=8 row of tables/invariance.tex.
% Strategy: argue the ORDER, not the picture.  a states the definition and b
% argues the addressing; both show a finished tree, which invites the reader to
% ask whether some other tree would have done.  This variant answers that by
% construction: once the leaf placement is fixed, no step below has a choice,
% so FBT(N) is not one canonical tree among many -- it is the only tree the
% rules can produce.  d makes the complementary argument from a counterexample.
% Colour: settled structure is spBlue, leaf slots spSky, and the level folded at
% each step is highlighted in spPurple so the eye lands on what changed; the
% closing root is the confirmed invariant (spGreen). The highlight is doubled by
% a heavier rule weight, so the step reads in greyscale too.
% Target: IEEE single column (3.5in = 8.9cm); drawn width approx 8.6cm.
\begin{tikzpicture}[
  x=1cm, y=1cm,
  inode/.style={draw=spBlue, circle, thick, inner sep=0pt, minimum size=4.2mm,
                font=\tiny, fill=spBlueF, text=spInk},
  new/.style={draw=spPurple!85!spInk, circle, line width=1.1pt, inner sep=0pt,
              minimum size=4.2mm, font=\tiny, fill=spPurple!16, text=spInk},
  lnode/.style={draw=spSky!80!spInk, rounded corners=0.8pt, thick, inner sep=0pt,
                minimum size=4.0mm, font=\tiny, fill=spSkyF, text=spInk},
  edg/.style={draw=spBlue!70, thin},
  nedg/.style={draw=spPurple!85!spInk, line width=0.9pt},
  cptn/.style={font=\tiny, align=left, anchor=north west, inner sep=0pt,
              text width=3.50cm, text=spInk},
  stp/.style={font=\tiny\bfseries, anchor=north west, inner sep=0pt, text=spBlue},
]
\def\p{0.53}          % leaf pitch
\def\lv{0.56}         % level spacing
\def\ox{0.30}         % left margin of the drawings
\def\tx{4.40}         % left margin of the captions

% ---------------------------------------------------------------------------
% One stage is drawn by \stage{<base y>}{<levels folded so far, 0..3>}
% Leaf i sits at x = \ox + i*\p; internal node centres are the midpoints of the
% leaf span they cover, which is what the heap indexing already fixes.
% ---------------------------------------------------------------------------
\newcommand{\stage}[3]{%
  % #1 = stage number, #2 = base y, #3 = levels folded
  \foreach \i in {0,...,7} {
    \pgfmathsetmacro{\xx}{\ox+\i*\p}
    \node[lnode] (S#1L\i) at (\xx,#2) {\pgfmathparse{int(\i+7)}\pgfmathresult};
  }
  \ifnum#3>0
    \foreach \n/\a/\b in {3/0/1, 4/2/3, 5/4/5, 6/6/7} {
      \pgfmathsetmacro{\xx}{\ox+(\a+\b)*\p/2}
      \pgfmathsetmacro{\yy}{#2+\lv}
      \ifnum#3=1 \node[new] (S#1N\n) at (\xx,\yy) {\n};
      \else      \node[inode] (S#1N\n) at (\xx,\yy) {\n};\fi
      \ifnum#3=1 \draw[nedg] (S#1N\n) -- (S#1L\a); \draw[nedg] (S#1N\n) -- (S#1L\b);
      \else      \draw[edg]  (S#1N\n) -- (S#1L\a); \draw[edg]  (S#1N\n) -- (S#1L\b);\fi
    }
  \fi
  \ifnum#3>1
    \foreach \n/\a/\b/\m in {1/3/4/1.5, 2/5/6/5.5} {
      \pgfmathsetmacro{\xx}{\ox+\m*\p}
      \pgfmathsetmacro{\yy}{#2+2*\lv}
      \ifnum#3=2 \node[new] (S#1N\n) at (\xx,\yy) {\n};
      \else      \node[inode] (S#1N\n) at (\xx,\yy) {\n};\fi
      \ifnum#3=2 \draw[nedg] (S#1N\n) -- (S#1N\a); \draw[nedg] (S#1N\n) -- (S#1N\b);
      \else      \draw[edg]  (S#1N\n) -- (S#1N\a); \draw[edg]  (S#1N\n) -- (S#1N\b);\fi
    }
  \fi
  \ifnum#3>2
    \pgfmathsetmacro{\xx}{\ox+3.5*\p}
    \pgfmathsetmacro{\yy}{#2+3*\lv}
    \node[new, draw=spGreen, fill=spGreenF] (S#1N0) at (\xx,\yy) {0};
    \draw[nedg, draw=spGreen] (S#1N0) -- (S#1N1); \draw[nedg, draw=spGreen] (S#1N0) -- (S#1N2);
  \fi
}

% ---- step 0 : the leaves, and only the leaves -------------------------------
\stage{0}{0.00}{0}
\node[stp] at (\tx,0.46) {Step 0 \ \ place the leaves};
\node[cptn]  at (\tx,0.18)
  {Slot $L[k]$ goes to heap index $k+(N{-}1)$, so $L[0..7]$ occupy $7..14$
   ascending. Nothing else is placed, and nothing else \emph{can} be: indices
   $0..N{-}2$ are not leaves.};

% ---- step 1 : fold level 3 ---------------------------------------------------
\stage{1}{1.62}{1}
\node[stp] at (\tx,2.66) {Step 1 \ \ fold once};
\node[cptn]  at (\tx,2.38)
  {Every leaf pair $(2n{+}1,2n{+}2)$ has one parent $n$. Leaves $7..14$ pair as
   $(7,8),(9,10),(11,12),(13,14)$ into $3,4,5,6$. No leaf may pair with a
   non-sibling: $\lfloor(n{-}1)/2\rfloor$ is a function.};

% ---- step 2 : fold level 2 ---------------------------------------------------
\stage{2}{3.62}{2}
\node[stp] at (\tx,5.18) {Step 2 \ \ fold again};
\node[cptn]  at (\tx,4.90)
  {The four new nodes are themselves a sibling set: $(3,4)\to1$, $(5,6)\to2$.
   The rule that produced them is the rule that consumes them, so the second
   level is forced by the first.};

% ---- step 3 : the root -------------------------------------------------------
\stage{3}{6.12}{3}
\node[stp] at (\tx,8.24) {Step 3 \ \ close at the root};
\node[cptn]  at (\tx,7.96)
  {$(1,2)\to0$ exhausts the index range: $2N{-}1=15$ nodes, all placed. The
   emission order is the post-order
   $\pi_8=\langle3,4,1,5,6,2,0\rangle$, and the root is
   $\Rcan(8,L)=$ \texttt{0x40E1948E978D4FDF} (Table~I).};
\node[font=\tiny, anchor=west, text=spGreen] at (2.63,7.80) {root};

% ---- the claim the sequence supports ----------------------------------------
\node[anchor=north west, font=\tiny, align=left, inner sep=0pt, text width=7.95cm,
      text=spInk] at (-0.20,-1.15)
  {At no step was there a choice. Fixing the leaf placement fixes every
   parent, because $\mathrm{parent}(n)=\lfloor(n{-}1)/2\rfloor$ is single-valued
   --- so $\FBT(N)$ is not a canonical representative selected from a family of
   admissible trees, it is the only tree the placement admits. That is what lets
   the ABI pin a bit-exact root without also pinning the topology, the mapping
   or the schedule: those choose \emph{where} each node runs, never \emph{which}
   pair it adds.};
\end{tikzpicture}
